%%%%%%%%%%%%%%%%%%%%%%%%%%%%%%%%%%%%%%%%%%%%%%%%%%%%%%%%%%%%%%%%%%%%%%%%%%%%
% CAPÍTULO 2: PROBLEMA Y PREGUNTAS DE INVESTIGACIÓN
%%%%%%%%%%%%%%%%%%%%%%%%%%%%%%%%%%%%%%%%%%%%%%%%%%%%%%%%%%%%%%%%%%%%%%%%%%%%

\chapter{Problema y preguntas de investigación}
\label{ch:problema}

\section{Formulación del problema}

El problema abordado por este TFM se concreta en la necesidad de diseñar, implementar y evaluar una plataforma de carpooling universitario que:
\begin{itemize}
    \item Facilitar la publicación y solicitud de viajes previamente planificados.
    \item Identificar compatibilidades entre conductores y pasajeros mediante mecanismos de \textit{matching}.
    \item Mantener una representación digital sincronizada del estado del sistema.
    \item Evaluar el comportamiento de la solución mediante simulación antes de su despliegue real.
\end{itemize}

Este problema se sitúa en la intersección de cuatro áreas: los Gemelos Digitales aplicados a movilidad, los sistemas de movilidad compartida P2P, los algoritmos de \textit{matching} conductor-pasajero y la simulación basada en agentes como instrumento de evaluación.

\section{Pregunta de investigación principal}

\textbf{Pregunta de investigación 1 (RQ1, del inglés \textit{Research Question}).} ¿Cómo puede una plataforma de carpooling universitaria apoyarse en un Gemelo Digital para representar, monitorizar y simular los desplazamientos compartidos y evaluar su impacto sobre el sistema de movilidad?

\section{Preguntas de investigación secundarias}

\textbf{RQ2.} ¿Qué entidades, atributos y eventos deben formar parte del modelo digital de usuarios, vehículos, viajes y rutas para que el Gemelo Digital resulte útil en el contexto del carpooling universitario?

\textbf{RQ3.} ¿Qué criterios de compatibilidad y qué mecanismos de \textit{matching} permiten identificar de forma eficiente oportunidades de viaje compartido en una comunidad universitaria con desplazamientos recurrentes?

\textbf{RQ4.} ¿Qué métricas e indicadores permiten evaluar el efecto del carpooling sobre el uso de vehículos y la ocupación en escenarios simulados, y cómo se comporta el sistema ante variaciones de demanda, horarios y capacidad?

\section{Relación entre preguntas, objetivos y categorías de la revisión}

La Tabla~\ref{tab:relacion-rq} relaciona cada pregunta de investigación con su objetivo asociado y las categorías de la revisión sistemática de la literatura (SLR, del inglés \textit{Systematic Literature Review}) utilizadas.

\begin{table}[H]
  \centering
  \begin{tabularx}{\textwidth}{Y{0.4154} Y{1.2692} Y{1.3154}}
    \hline
    \textbf{Pregunta} & \textbf{Objetivo asociado} & \textbf{Categorías de la SLR relacionadas} \\
    \hline
    RQ1 & Diseñar la arquitectura de la plataforma y del Gemelo Digital & A (Fundamentos y arquitectura DT), B (DT + movilidad compartida) \\
    RQ2 & Definir el modelo digital de usuarios, vehículos, viajes y rutas & A, C (Datos, integración y sincronización) \\
    RQ3 & Diseñar e implementar la lógica de \textit{matching} & E (\textit{Matching}, rutas, optimización) \\
    RQ4 & Definir métricas y evaluar el sistema mediante simulación & D (Simulación y escenarios), E (Evaluación) \\
    \hline
  \end{tabularx}
  \caption{Relación entre preguntas, objetivos y categorías de la SLR.}
  \label{tab:relacion-rq}
\end{table}
